\subsection{Smoothness and positivity of the density}\label{sec:density}

The smoothness conclusion and its deduction from saturated consistency
appear in Chen~\cite[Theorem~1.3 and Proposition~4.6]{Chen2026}.
We include the argument to show how Theorem~\ref{thm:full-support}
determines the regularity of the minimizing order parameter.
The polynomial-moment differentiation follows Auffinger and
Chen~\cite[Lemma~2 and the proof of Theorem~2(ii)]{ACProperties2015},
in its zero-temperature form in Chen~\cite[Appendix~B.6]{Chen2026};
the Hessian semimartingale identity is also given by Chen, Handschy,
and Lerman~\cite[Lemma~3]{CHL2018}. We then use the coercive inequality
of Section~\ref{cert:section} to prove that the density is strictly
positive at every positive time below one.

\begin{corollary}[Smooth positive density of the minimizing measure]\label{cor:smooth-measure}
Let $\gamma\in\mathcal U$ minimize $\mathcal P$, and let $X$ be its
controlled diffusion. Then $\gamma(0)=0$ and $\gamma$ is smooth on
$[0,1)$, with derivatives at zero understood from the right. Moreover,
\begin{equation}\label{eq:gamma-moment-ratio}
 \gamma(t)=
 \frac{\E[\Phi_{\gamma,xxx}(t,X_t)^2]}
      {2\E[\Phi_{\gamma,xx}(t,X_t)^3]},
 \qquad 0\leq t<1.
\end{equation}
Consequently $\mu_\gamma(\dd t)=\gamma'(t)\,\dd t$ on $[0,1)$,
where $\gamma'$ is smooth and nonnegative. In fact, for $0<t<1$,
\begin{equation}\label{eq:positive-density-bound}
 2\gamma'(t)\E[\Phi_{\gamma,xx}(t,X_t)^3]
 \geq \frac{\gamma(t)^2}{1000}
                \E[\Phi_{\gamma,xx}(t,X_t)^4]>0.
\end{equation}
In particular $\gamma'(t)>0$ on $(0,1)$.
\end{corollary}

\begin{proof}
Write $f_k=\partial_x^k\Phi_\gamma$ for $k\geq1$, and evaluate these
functions along $X$ unless otherwise indicated. All spatial derivatives
are bounded on each compact time strip by Lemma~\ref{lem:regularity}.
Differentiating the PDE and applying It\^o's formula gives
\begin{equation}\label{eq:spatial-derivative-martingales}
 \dd f_k(t,X_t)
 =-\frac{\gamma(t)}2\sum_{j=1}^{k-1}\binom{k}{j}
       f_{j+1}(t,X_t)f_{k-j+1}(t,X_t)\,\dd t
       +f_{k+1}(t,X_t)\,\dd B_t.
\end{equation}
The sum is empty for $k=1$. For general integrable $\gamma$, this
identity follows by the finite-step approximation and It\^o isometry
used in Lemma~\ref{lem:diffusion}. On a fixed strip the spatial
factors converge uniformly along the coupled paths, the time coefficients
converge in $L^1$, and the stochastic integrals converge in $L^2$.

For $a=f_2$ and $b=f_3$, the identity is
$\dd a=-\gamma a^2\,\dd t+b\,\dd B_t$. Hence
\[
 \frac{\dd}{\dd t}\E a(t,X_t)^2
 =\E b(t,X_t)^2-2\gamma(t)\E a(t,X_t)^3
 \quad\hbox{for almost every }t<1.
\]
The left side is zero by Theorem~\ref{thm:full-support}. Convexity gives
$a\geq0$, and the same theorem gives $\E a^2=1$, so
$\E a^3\geq(\E a^2)^{3/2}=1$. The ratio on the right of
\eqref{eq:gamma-moment-ratio} is therefore well defined and continuous
on every compact time strip. Its almost-everywhere equality with
$\gamma$ extends to every time by right continuity. In particular
$\gamma$ is continuous. At time zero, evenness gives $b(0,0)=0$,
so the formula yields $\gamma(0)=0$.

To bootstrap regularity without assuming time derivatives of $\gamma$,
consider all polynomial moments of the spatial derivatives at once.
For a polynomial $P=P(y_1,\ldots,y_m)$, define
\[
 M_P(t)=\E P(f_1(t,X_t),\ldots,f_m(t,X_t)).
\]
In formal variables $y_1,y_2,\ldots$, set
\[
 \begin{aligned}
 D_k(y)&=-\frac12\sum_{j=1}^{k-1}\binom{k}{j}
                           y_{j+1}y_{k-j+1},\\
 Q_P(y)&=\frac12\sum_{k,l=1}^m
              (\partial_k\partial_lP)(y)y_{k+1}y_{l+1},\\
 R_P(y)&=\sum_{k=1}^m(\partial_kP)(y)D_k(y).
 \end{aligned}
\]
These are polynomials in finitely many variables, with $D_1=0$.
It\^o's product rule and~\eqref{eq:spatial-derivative-martingales} give
\begin{equation}\label{eq:polynomial-moment-bootstrap}
 M_P(t)-M_P(0)=\int_0^t
       \bigl(M_{Q_P}(s)+\gamma(s)M_{R_P}(s)\bigr)\,\dd s.
\end{equation}
Every such moment is finite and continuous on each compact strip.
Since $\gamma$ is continuous, \eqref{eq:polynomial-moment-bootstrap}
makes every $M_P$ continuously differentiable; the ratio
\eqref{eq:gamma-moment-ratio} then makes $\gamma$ continuously
differentiable. More generally, if $\gamma$ and all these moments are
$C^k$, the integrand in \eqref{eq:polynomial-moment-bootstrap} is $C^k$
for every $P$, making all moments $C^{k+1}$. The same ratio makes
$\gamma$ $C^{k+1}$. Induction proves smoothness, including right
derivatives at zero.

Smoothness and $\gamma(0)=0$ identify its Stieltjes measure with
$\gamma'(t)\,\dd t$, and monotonicity gives $\gamma'\geq0$.
It remains to prove strict positivity away from zero. Set $d=f_4$.
The polynomial-moment identities give
\[
 \frac{\dd}{\dd t}\E b^2=\E d^2-6\gamma\E(ab^2),
 \qquad
 \frac{\dd}{\dd t}\E a^3=3\E(ab^2)-3\gamma\E a^4.
\]
Differentiating $\E b^2=2\gamma\E a^3$ therefore yields
\begin{equation}\label{eq:density-curvature-identity}
 2\gamma'\E a^3
  =\E d^2-12\gamma\E(ab^2)+6\gamma^2\E a^4.
\end{equation}
All factors in this identity are evaluated at $(t,X_t)$.

Fix $t\in(0,1)$ and put $c=\gamma(t)$. Full support implies
$c>0$. We next apply the integrated coercivity estimate at this fixed
time, although $\gamma$ need not be constant on a neighborhood of $t$.
Choose $\delta_n\downarrow0$ with
$[t-\delta_n,t+\delta_n]\subset(0,1)$. Replace $\gamma$ on
$[t-\delta_n,t+\delta_n)$ by the value $c$. This preserves
monotonicity, since $\gamma(s)\leq c$ for $s<t$ and
$\gamma(s)\geq c$ for $s>t$, and the change tends to zero in $L^1$.
Outside that interval, truncate at a height greater than $c$, impose
a positive lower cutoff less than $c$, and approximate separately on
the left and right by nondecreasing finite step functions with values
at most $c$ and at least $c$, respectively. Let the truncation height
tend to infinity, the lower cutoff tend to zero, and the discretization
error tend to zero. The resulting positive bounded nondecreasing
coefficients $\gamma_n$ converge to $\gamma$ in $L^1$ and equal $c$
on a neighborhood of $t$ for every $n$.

Use terminal data $M_n^{-1}\log(2\cosh(M_nx))$, with
$M_n\geq\|\gamma_n\|_\infty$ and $M_n\to\infty$, and let
$a_n,b_n,d_n$ and $X^n$ denote the corresponding spatial derivatives
and diffusion. At time $t$, the scaled derivatives in
Proposition~\ref{cert:coercivity} are
$v_n=ca_n$, $w_n=cb_n$, and $z_n=cd_n$. Thus that proposition and
Lemmas~\ref{shape:backward}--\ref{shape:ibp} give
\[
 \E[d_n^2-12c a_n b_n^2+6c^2a_n^4](t,X_t^n)
       \geq\frac{c^2}{1000}\E a_n(t,X_t^n)^4.
\]
The derivative bounds and convergence on a fixed compact strip from
Lemmas~\ref{lem:regularity} and~\ref{lem:diffusion} do not depend on
the widths $\delta_n$ of the artificial plateaus. They therefore
allow passage to the limit in every expectation in the last display.
Combine the limiting inequality with
\eqref{eq:density-curvature-identity} and $c=\gamma(t)$ to obtain
\eqref{eq:positive-density-bound}. Its right side is strictly positive
because $c>0$ and $\E a^4\geq(\E a^2)^2=1$. The proof asserts only
$\gamma'(0)\geq0$ at the left endpoint.
\end{proof}
